% --- ARCHIVO: Anexos_Conceptos/anexo_cap4.tex ---

\chapter{Conceptos técnicos y regulatorios del Capítulo 4}
\label{anexo:conceptos_cap4}

\section{Procedimiento de Operación 1.6 (P.O. 1.6)}
\label{anexo:po16}
El Procedimiento de Operación 1.6 es el protocolo de emergencia del sistema eléctrico español que establece los planes de salvaguarda y reposición del suministro ante incidentes críticos. Dictamina las estrategias de fragmentación topológica de la red en islas eléctricas independientes, las rutas de energización preferentes y el protocolo de priorización de arranque de las instalaciones de generación para restaurar el sistema tras un cero de tensión parcial o total.

\section{Arranque autónomo (Black Start)}
\label{anexo:black_start}
La capacidad de \textit{Black Start} (arranque autónomo o arranque en negro) es el servicio de ajuste por el cual ciertas instalaciones de generación pueden arrancar y comenzar a inyectar energía a la red sin necesidad de recibir tensión eléctrica externa. Esta capacidad, tradicionalmente provista por centrales hidroeléctricas con válvulas de apertura mecánica o por grupos térmicos con generadores diésel de emergencia, es el requisito primario para iniciar la re-energización (estrategia \textit{Bottom-Up}) de un sistema eléctrico tras un colapso total.

\section{Potencia de cortocircuito ($S_{sc}$)}
\label{anexo:potencia_cortocircuito}
La potencia de cortocircuito ($S_{sc}$) en un nudo de la red es una medida de su "fortaleza" electromagnética. Representa la cantidad de corriente aparente que fluiría hacia ese nudo en caso de producirse un cortocircuito franco trifásico. Una elevada potencia de cortocircuito, típicamente aportada por los grandes generadores síncronos, implica que la tensión en ese nudo es muy robusta y resiliente, sufriendo variaciones mínimas ante perturbaciones, conexiones de cargas bruscas o maniobras en la red.

\section{Relés de comprobación de sincronismo (Synchro-check)}
\label{anexo:synchro_check}
El relé \textit{synchro-check} (función 25 ANSI) es un dispositivo de protección empleado en las maniobras de acoplamiento de sistemas eléctricos separados (islas). Su función es supervisar continuamente que la tensión, la frecuencia y el ángulo de fase a ambos lados de un interruptor abierto se encuentran dentro de unos márgenes de tolerancia preestablecidos. Solo cuando estas tres variables coinciden, el relé autoriza el cierre del interruptor, garantizando que las islas se sincronicen de forma segura sin provocar un impacto transitorio severo.

\section{Centros de Coordinación Regional (RCC)}
\label{anexo:rcc}
Los Centros de Coordinación Regional (RCC, \textit{Regional Coordination Centres}) son entidades supranacionales establecidas por la normativa europea para facilitar la cooperación operativa entre los distintos Gestores de Redes de Transporte (TSO). Su función no es operar los interruptores en tiempo real, sino realizar análisis de seguridad coordinados (CSA), calcular capacidades de intercambio transfronterizo y proponer acciones preventivas y paliativas para asegurar el cumplimiento del Criterio N-1 a nivel regional (ej. Coreso para la región del Suroeste de Europa).

\section{Reserva de Restauración de Frecuencia Automática (aFRR)}
\label{anexo:afrr}
La aFRR (\textit{automatic Frequency Restoration Reserve}), históricamente conocida como regulación secundaria, es un servicio de ajuste del sistema que se activa automáticamente tras una desviación de frecuencia. Está controlado directamente por el AGC (\textit{Automatic Generation Control}) del Operador del Sistema y su objetivo es devolver progresivamente la frecuencia a su valor nominal (50 Hz) y restituir los flujos en las interconexiones a sus programas pactados, liberando así la reserva primaria (FCR) previamente activada.

\section{Centro de Control de Energías Renovables (CECRE)}
\label{anexo:cecre}
El CECRE es el centro de supervisión y control dependiente de Red Eléctrica de España, pionero a nivel mundial, encargado de monitorizar y gobernar la inyección de la generación basada en energías renovables (régimen RCR). Su función es garantizar que la integración masiva de fuentes de generación de naturaleza intermitente (eólica y solar) se realice de forma segura, teniendo potestad para enviar consignas de reducción de potencia (curtailment) cuando la seguridad del sistema lo requiere.